%==============================================================
%  lecons/analyse/lagrange.tex — Interpolation de Lagrange
%==============================================================

\lecontitle{Les polynômes de Lagrange}{Analyse réelle — Interpolation polynomiale}

%=== NOTATION LOCALE =========================================
\newcommand{\Lbase}[1]{L_{#1}}
\newcommand{\Rn}{\mathbb{R}_n[X]}

%=============================================================
%  § 1 — DÉFINITION ET RÉSULTATS FONDAMENTAUX
%=============================================================
\secnum{navyblue}{1}{Définition et résultats fondamentaux}

\begin{tcolorbox}[thmbox]
  \textbf{Définition.}\enspace\index{polynômes interpolateurs de Lagrange}%
  \index{interpolation polynomiale}
  Soit $x_0, x_1, \ldots, x_n$ des réels \emph{deux à deux distincts}. On
  appelle \emph{polynômes de base de Lagrange} associés à ces nœuds les
  $n+1$ polynômes
  \[
    \Lbase{i}(x) \;=\; \prod_{\substack{0 \leq j \leq n \\ j \neq i}}
    \frac{x - x_j}{x_i - x_j}, \qquad 0 \leq i \leq n.
  \]
  Chaque $\Lbase{i}$ est de degré exactement $n$ et vérifie la propriété
  fondamentale
  \[
    \Lbase{i}(x_k) = \delta_{ik} =
    \begin{cases} 1 & \text{si } k = i, \\ 0 & \text{si } k \neq i, \end{cases}
  \]
  où $\delta_{ik}$ désigne le symbole de Kronecker.

  \medskip
  \textbf{Théorème (existence, unicité, erreur).}
  \begin{itemize}
    \item \textit{Existence et unicité.}\enspace Pour toute famille de
          couples ${(x_i, y_i)}_{0 \leq i \leq n}$ avec les $x_i$ distincts,
          il existe un \emph{unique} polynôme $P \in \Rn$ tel que
          $P(x_i) = y_i$ pour tout $i$, donné par
          \[
            P(x) \;=\; \sum_{i=0}^{n} y_i\, \Lbase{i}(x).
          \]
    \item \textit{Base.}\enspace La famille ${(\Lbase{i})}_{0 \leq i \leq n}$
          est une base de $\Rn$, duale de la famille des formes
          d'évaluation $\varphi_k : Q \mapsto Q(x_k)$.
    \item \textit{Formule d'erreur.}\enspace Si $f$ est de classe
          $\mathcal{C}^{n+1}$ sur un intervalle $I$ contenant les $x_i$ et
          $x$, et si $P$ interpole $f$ aux nœuds ($y_i = f(x_i)$), alors il
          existe $\xi \in I$ tel que
          \[
            f(x) - P(x) \;=\; \frac{f^{(n+1)}(\xi)}{(n+1)!}
            \prod_{j=0}^{n} (x - x_j).
          \]
    \item \textit{Lien avec Vandermonde.}\enspace\index{Vandermonde!matrice de}
          Dans la base canonique $(1, X, \ldots, X^n)$, la recherche de $P$
          équivaut à résoudre un système linéaire dont la matrice est la
          matrice de Vandermonde $V = {\bigl(x_i^{\,j}\bigr)}_{0 \leq i,j \leq n}$,
          inversible car $\det V = \prod_{i<j}(x_j - x_i) \neq 0$.
  \end{itemize}
\end{tcolorbox}

\begin{significationintuitive}
Interpoler, c'est faire passer une courbe polynomiale par des points
imposés du plan. Le théorème affirme qu'avec $n+1$ abscisses distinctes,
il existe une et une seule telle courbe de degré $\leq n$. Les $\Lbase{i}$
sont les « briques élémentaires » : $\Lbase{i}$ vaut $1$ au nœud $x_i$ et
$0$ aux autres, de sorte que pondérer par les $y_i$ reconstitue exactement
les ordonnées voulues. Algébriquement, $(\Lbase{i})$ est la base duale des
formes d'évaluation aux nœuds : connaître un polynôme de $\Rn$ équivaut à
connaître ses valeurs en $n+1$ points.
\end{significationintuitive}

\decsep{}

%=============================================================
%  § 2 — DÉMONSTRATION
%=============================================================
\secnum{navyblue}{2}{Démonstration}

\begin{ideepreuve}
L'existence est constructive : le polynôme $\sum y_i \Lbase{i}$ convient
grâce à la propriété $\Lbase{i}(x_k)=\delta_{ik}$. L'unicité repose sur un
argument de degré : une différence de deux interpolateurs est de degré
$\leq n$ et possède $n+1$ racines, donc est nulle. La formule d'erreur
s'obtient en figeant $x$ et en appliquant le théorème de Rolle de façon
itérée à une fonction auxiliaire bien choisie.
\end{ideepreuve}

\vspace{10pt}
\rubriq{Preuve}

\begin{mdframed}[style=proofbar]
  \textit{Étape 1 — Les $\Lbase{i}$ et la propriété
  $\Lbase{i}(x_k)=\delta_{ik}$.}
  Par construction, $\Lbase{i}$ est un produit de $n$ facteurs affines
  $\tfrac{x-x_j}{x_i-x_j}$ ($j \neq i$), donc $\deg \Lbase{i} = n$ et le
  dénominateur $\prod_{j\neq i}(x_i-x_j)$ est non nul (nœuds distincts).
  Si $k \neq i$, le facteur d'indice $j=k$ s'annule en $x=x_k$, donc
  $\Lbase{i}(x_k)=0$. Si $k=i$, chaque facteur vaut
  $\tfrac{x_i-x_j}{x_i-x_j}=1$, donc $\Lbase{i}(x_i)=1$.

  \medskip
  \textit{Étape 2 — Existence.}
  Posons $P = \sum_{i=0}^{n} y_i\,\Lbase{i}$. C'est un polynôme de degré
  $\leq n$ comme combinaison linéaire des $\Lbase{i}$. Pour chaque nœud
  $x_k$ :
  \[
    P(x_k) = \sum_{i=0}^{n} y_i\,\Lbase{i}(x_k)
           = \sum_{i=0}^{n} y_i\,\delta_{ik} = y_k.
  \]
  Donc $P$ interpole les données : l'existence est acquise.

  \medskip
  \textit{Étape 3 — Unicité.}
  Soit $P_1, P_2 \in \Rn$ deux interpolateurs des mêmes données. Le
  polynôme $D = P_1 - P_2$ est de degré $\leq n$ et s'annule aux $n+1$
  points distincts $x_0, \ldots, x_n$. Un polynôme non nul de degré
  $\leq n$ a au plus $n$ racines ; ayant $n+1$ racines, $D$ est le
  polynôme nul, d'où $P_1 = P_2$. En particulier, $(\Lbase{i})$ engendre
  $\Rn$ (étape~2 avec des $y_i$ arbitraires) et compte $n+1 = \dim \Rn$
  éléments : c'est une base.

  \medskip
  \textit{Étape 4 — Formule d'erreur.}
  Fixons $x \in I$. Si $x = x_k$ pour un certain $k$, les deux membres
  sont nuls. Supposons donc $x \notin \{x_0,\ldots,x_n\}$ et posons
  $\omega(t) = \prod_{j=0}^{n}(t-x_j)$, puis
  \[
    K = \frac{f(x) - P(x)}{\omega(x)}, \qquad
    g(t) = f(t) - P(t) - K\,\omega(t).
  \]
  La constante $K$ est choisie de sorte que $g(x)=0$. De plus
  $g(x_k) = f(x_k) - P(x_k) - K\,\omega(x_k) = 0$ pour tout $k$, car
  $P(x_k)=f(x_k)$ et $\omega(x_k)=0$. Ainsi $g$ s'annule en les $n+2$
  points distincts $x, x_0, \ldots, x_n$. Comme $f \in \mathcal{C}^{n+1}$
  (et $P,\omega$ polynomiaux), $g$ est de classe $\mathcal{C}^{n+1}$.

  En appliquant le théorème de Rolle entre chaque paire de zéros
  consécutifs, $g'$ admet au moins $n+1$ zéros ; en réitérant, $g''$ en a
  au moins $n$, \ldots, et $g^{(n+1)}$ admet au moins un zéro $\xi \in I$.
  Or $P^{(n+1)} = 0$ ($\deg P \leq n$) et $\omega^{(n+1)} = (n+1)!$
  ($\omega$ unitaire de degré $n+1$), donc
  \[
    g^{(n+1)}(t) = f^{(n+1)}(t) - K\,(n+1)!.
  \]
  En $t=\xi$ : $0 = f^{(n+1)}(\xi) - K\,(n+1)!$, d'où
  $K = \dfrac{f^{(n+1)}(\xi)}{(n+1)!}$. En reportant dans la définition de
  $K$ :
  \[
    f(x) - P(x) = K\,\omega(x)
    = \frac{f^{(n+1)}(\xi)}{(n+1)!} \prod_{j=0}^{n}(x-x_j).
  \]
  \hfill$\blacksquare$
\end{mdframed}

\decsep{}

%=============================================================
%  § 3 — EXEMPLES, CONTRE-EXEMPLES, EXERCICES
%=============================================================
\secnum{exgreen}{3}{Exemples, contre-exemples et exercices}

\rubriq{Exemples}

\begin{mdframed}[style=exbar]
  \textbf{1. Interpolation par trois points.}\enspace
  Cherchons $P \in \mathbb{R}_2[X]$ tel que $P(0)=1$, $P(1)=3$, $P(2)=2$.
  Avec $x_0=0$, $x_1=1$, $x_2=2$ :
  \[
    \Lbase{0}(x)=\tfrac{(x-1)(x-2)}{2},\quad
    \Lbase{1}(x)=-x(x-2),\quad
    \Lbase{2}(x)=\tfrac{x(x-1)}{2}.
  \]
  D'où
  \[
    P = 1\cdot\Lbase{0} + 3\cdot\Lbase{1} + 2\cdot\Lbase{2}
      = -\tfrac{3}{2}x^2 + \tfrac{7}{2}x + 1.
  \]

  \smallskip\noindent
  \textbf{2. Base de $\mathbb{R}_2[X]$ et lien Vandermonde.}\enspace
  Les trois polynômes $\Lbase{0},\Lbase{1},\Lbase{2}$ forment une base de
  $\mathbb{R}_2[X]$. La matrice de changement de base depuis
  $(1,X,X^2)$ est l'inverse de la matrice de Vandermonde
  \[
    V = \begin{pmatrix} 1 & 0 & 0 \\ 1 & 1 & 1 \\ 1 & 2 & 4 \end{pmatrix},
    \qquad \det V = (1-0)(2-0)(2-1) = 2 \neq 0.
  \]

  \smallskip\noindent
  \textbf{3. Forme de Newton.}\enspace%
  \index{différences divisées}
  Pour un calcul \emph{incrémental} (ajout d'un nœud sans tout recommencer),
  on écrit le \emph{même} polynôme $P$ dans la base de Newton, ses
  coefficients étant les différences divisées de $f$ : cf.\ la leçon
  suivante, \emph{La forme de Newton et les différences divisées}.
\end{mdframed}

\vspace{6pt}
\rubriq{Contre-exemples et pièges}

\begin{mdframed}[style=cexbar]
  \textbf{Phénomène de Runge.}\enspace\index{phénomène de Runge}
  Interpoler $f(x)=\dfrac{1}{1+25x^2}$ sur $[-1,1]$ aux $n+1$ nœuds
  \emph{équirépartis} produit une suite $(P_n)$ qui \emph{diverge} :
  $\max_{[-1,1]}|f-P_n| \to +\infty$, avec des oscillations explosant près
  des bords. Augmenter le degré aggrave l'erreur, contrairement à
  l'intuition.

  \smallskip\noindent
  \textbf{Remède : nœuds de Tchebychev.}\enspace\index{nœuds de Tchebychev}
  La formule d'erreur fait intervenir $\|\omega\|_\infty$ avec
  $\omega = \prod (x-x_j)$. Choisir pour $x_j$ les racines de $T_{n+1}$
  (nœuds de Tchebychev, cf.\ leçon \emph{Polynômes de Tchebychev})
  minimise $\|\omega\|_\infty$ et restaure la convergence pour les
  fonctions régulières.

  \smallskip\noindent
  \textbf{Nœuds distincts obligatoires.}\enspace
  Si deux $x_i$ coïncident, la matrice de Vandermonde est singulière et le
  problème n'a plus de solution unique (il faudrait imposer aussi des
  dérivées : interpolation d'Hermite).

  \smallskip\noindent
  \textbf{Instabilité numérique.}\enspace
  En haut degré, l'évaluation directe de $\sum y_i \Lbase{i}$ est sensible
  aux erreurs d'arrondi ; on lui préfère la \emph{forme barycentrique} de
  Lagrange, numériquement stable : en posant
  $w_i = 1/\prod_{j\neq i}(x_i-x_j)$,
  \[
    P(x) = \frac{\displaystyle\sum_{i=0}^{n} \frac{w_i\,y_i}{x-x_i}}
                {\displaystyle\sum_{i=0}^{n} \frac{w_i}{x-x_i}}
    \qquad \bigl(x \notin \{x_0,\ldots,x_n\}\bigr),
  \]
  le dénominateur provenant de l'identité $\sum_i \Lbase{i} = 1$
  (exercice~2).
\end{mdframed}

\vspace{6pt}
\exolabel

\begin{mdframed}[style=exobar]
  \begin{enumerate}
    \item Déterminer le polynôme $P \in \mathbb{R}_2[X]$ tel que
          $P(-1)=2$, $P(0)=0$, $P(1)=2$. \textit{Indication :} écrire les
          trois $\Lbase{i}$, ou remarquer la parité.

    \item Montrer que $\sum_{i=0}^{n} \Lbase{i}(x) = 1$ pour tout $x$.
          \textit{Indication :} c'est l'interpolateur de la fonction
          constante $1$, dont l'unicité conclut.

    \item Établir l'identité $\sum_{i=0}^{n} \Lbase{i}'(x_k) = 0$ pour tout
          $k$, en dérivant la relation de la question~2.

    \item Soit $f(x)=\sin x$ interpolée aux nœuds $0, \tfrac{\pi}{2}, \pi$.
          Majorer $\max_{[0,\pi]}|f-P|$ à l'aide de la formule d'erreur
          (on a $|f^{(3)}| \leq 1$).

    \item (Différences divisées, cf.\ la leçon suivante) Construire la
          table des différences divisées de $f$ aux nœuds $0,1,2,3$ avec
          $f(0)=1, f(1)=0, f(2)=5, f(3)=22$, puis écrire $P$ sous forme de
          Newton et vérifier que $P(x)=x^3-2x+1$.

    \item (Vandermonde) Démontrer la formule
          $\det V = \prod_{0\leq i<j\leq n}(x_j-x_i)$ par récurrence sur
          $n$, en utilisant des opérations élémentaires sur les colonnes.
  \end{enumerate}
\end{mdframed}

\leconfooter{J.-L.\ Lagrange (1795) ; E.\ Waring (1779)}
